\thispagestyle{empty}
\section*{Avant-propos}\label{avant-propos}
\markboth{Avant-propos}{Avant-propos}

Ce document de référence du protocole APRS marque la maturité du bébé de
WB4APR. Parti d'un concept simple --- un moyen de suivre la position
d'objets mobiles par packet radio --- les programmes utilisant le
protocole APRS sont devenus peut-être l'application packet radio la plus
populaire aujourd'hui. C'est aussi devenue l'une des plus complexes ; du
concept initial est né, et continue à croître, un système de
communication tactique d'une capacité considérable. Comme beaucoup de
projets amateurs, le protocole APRS a été conçu au fil de son
implémentation, et beaucoup de ses subtilités n'ont jamais été
documentées.

Jusqu'à maintenant. Cette spécification définit le protocole APRS sur
les ondes avec une précision et une clarté qui en font un modèle pour
les efforts futurs. Le travail accompli par les membres de l'APRS
Working Group, ainsi que par l'éditeur technique Ian Wade, G3NRW, doit
être reconnu comme une contribution majeure à l'art du packet radio.
Avec ce document disponible, aucun développeur n'a plus d'excuse pour
implémenter incorrectement le protocole APRS.

En tant que membre de l'APRS Working Group dont le rôle était
principalement celui d'observateur, j'ai été fasciné par les échanges
entre les auteurs APRS et l'éditeur technique au fur et à mesure que la
spécification prenait forme. Mettre sur papier des détails qui
existaient auparavant uniquement dans l'esprit des auteurs a mis au jour
des ambiguïtés, des conséquences non envisagées, et même des erreurs
dans ce que les auteurs croyaient savoir. Les discussions qui ont suivi
chaque brouillon, et les questions que Ian a posées pour lever les
incertitudes, ont donné à chacun une meilleure compréhension du
protocole. Je suis sûr que ce processus a déjà contribué à une meilleure
interopérabilité entre les applications APRS existantes.

Tous ceux qui ont suivi le développement de la spécification, depuis sa
première mention en avril 1999 jusqu'à la publication de cette Version
1.0 en août 2000, savent que le processus a pris beaucoup plus de temps
que prévu. Dans le même temps, ils ont vu le brouillon se transformer
d'un squelette en un livre substantiel de plus de 110 pages. Avec la
spécification maintenant en main, je pense que nous pouvons tous dire
que l'attente en valait la peine. Félicitations à l'APRS Working Group
et, en particulier, à G3NRW, pour une contribution majeure à la
littérature du packet radio.

\emph{John Ackermann, N8UR --- TAPR Vice President et Administrative
Chair du APRS Working Group --- août 2000}
